% =========================================================
% Totally Bounded Sets
% Source: volume-iii/analysis/metric-spaces/notes/metric-spaces/
%
% Dependencies:
%   notes-metric-space.tex           — def:metric, def:open-ball
%   notes-properties-metric-space.tex — def:bounded-set, def:diameter
%   notes-completeness.tex           — completeness (for Generalized Heine-Borel)
% =========================================================

% ---------------------------------------------------------
\subsubsection{Totally Bounded Sets}
\label{sec:totally-bounded}
% ---------------------------------------------------------

% ---- Toolkit ----
\begin{tcolorbox}[colback=gray!6, colframe=gray!40, arc=2pt,
  left=6pt, right=6pt, top=4pt, bottom=4pt,
  title={\small\textbf{Totally Bounded Sets --- Quick Reference}},
  fonttitle=\small\bfseries]
\small
\begin{tabular}{l l l l}
\toprule
\textbf{Label} & \textbf{Statement} & \textbf{Proof method} & \textbf{Detail} \\
\midrule
D\ref{def:totally-bounded}
  & $A$ totally bounded: $\forall\varepsilon>0$, finitely many $\varepsilon$-balls cover $A$
  & Definition
  & \hyperref[def:totally-bounded]{$\downarrow$} \\[4pt]
P\ref{prop:totally-bounded-implies-bounded}
  & Totally bounded $\Rightarrow$ bounded
  & One ball covers all
  & \hyperref[prop:totally-bounded-implies-bounded]{$\downarrow$} \\[4pt]
P\ref{prop:bounded-not-totally-bounded}
  & Bounded $\not\Rightarrow$ totally bounded (infinite discrete)
  & Counterexample
  & \hyperref[prop:bounded-not-totally-bounded]{$\downarrow$} \\[4pt]
P\ref{prop:totally-bounded-Rn}
  & Bounded $\iff$ totally bounded in $\mathbb{R}^n$
  & Finite grid
  & \hyperref[prop:totally-bounded-Rn]{$\downarrow$} \\[4pt]
P\ref{prop:totally-bounded-Cauchy}
  & Every sequence in a totally bounded set has a Cauchy subsequence
  & Diagonal argument
  & \hyperref[prop:totally-bounded-Cauchy]{$\downarrow$} \\[4pt]
T\ref{thm:gen-heine-borel}
  & Compact $\iff$ complete $+$ totally bounded (Generalized Heine-Borel)
  & Statement (proof in compactness chapter)
  & \hyperref[thm:gen-heine-borel]{$\downarrow$} \\[4pt]
\bottomrule
\end{tabular}
\end{tcolorbox}

% ---------------------------------------------------------
\paragraph{Motivation.}
We know that \emph{bounded} sets are those contained in some ball of
finite radius.
But boundedness is a weak condition in general metric spaces.
The unit ball in $\ell^2$ (the space of square-summable sequences)
is bounded, but it is not compact --- you can find sequences
in it with no convergent subsequence.
Something extra is needed.

The missing ingredient is \emph{total boundedness}: the requirement
that, no matter how small you make $\varepsilon$, you can still cover
the set with \emph{finitely many} $\varepsilon$-balls.
In $\mathbb{R}^n$ this turns out to be equivalent to ordinary boundedness,
which is why the distinction never appeared in single-variable calculus.
In general metric spaces the two concepts diverge, and total boundedness
is the correct one for the theory of compactness.

% ---- Definition ----
\begin{tcolorbox}[colback=propbox, colframe=propborder, arc=2pt,
  left=6pt, right=6pt, top=4pt, bottom=4pt,
  title={\small\textbf{Definition (Totally Bounded Set)}},
  fonttitle=\small\bfseries]
\begin{definition}[Totally Bounded]
\label{def:totally-bounded}
Let $(X, d)$ be a metric space and $A \subseteq X$.
We say $A$ is \emph{totally bounded} if for every $\varepsilon > 0$ there
exist finitely many points $x_1, \ldots, x_n \in X$ such that
\[
A \;\subseteq\; \bigcup_{k=1}^n B_\varepsilon(x_k).
\]
The finite collection $\{x_1, \ldots, x_n\}$ is called an
\emph{$\varepsilon$-net} for $A$.
\end{definition}
\end{tcolorbox}

\begin{remark*}[Fully quantified form]
\[
A \text{ totally bounded}
\;\iff\;
\forall\,\varepsilon > 0,\;
\exists\, n \in \mathbb{N},\;
\exists\, x_1, \ldots, x_n \in X,\quad
A \subseteq \bigcup_{k=1}^n B_\varepsilon(x_k).
\]
\end{remark*}

\begin{remark*}[Intuition: covering by smaller and smaller finite sets]
Bounded says: $A$ fits inside \emph{one} ball of some fixed radius.
Totally bounded says: for any $\varepsilon$ you name --- no matter how
small --- $A$ fits inside \emph{finitely many} balls of radius $\varepsilon$.

The key word is ``finitely many.''
As $\varepsilon \to 0$, the balls shrink, but you must always be able to
cover $A$ with a finite number of them.
If $A$ is ``spread out'' in infinitely many directions that cannot be
captured by a finite set of small balls, then $A$ is bounded but not
totally bounded.
\end{remark*}

\begin{remark*}[Centers need not lie in $A$]
The definition does not require $x_1, \ldots, x_n \in A$.
But if you prefer, you can always choose centers in $A$:
if $A \subseteq \bigcup_k B_\varepsilon(x_k)$, pick any $a_k \in A \cap B_\varepsilon(x_k)$;
then $A \subseteq \bigcup_k B_{2\varepsilon}(a_k)$.
So totally bounded is equivalent to: for every $\varepsilon > 0$ there
is a finite $\varepsilon$-net \emph{inside} $A$ (up to doubling $\varepsilon$).
\end{remark*}

% ---- Totally bounded implies bounded ----

\begin{proposition}[Totally bounded implies bounded]
\label{prop:totally-bounded-implies-bounded}
If $A \subseteq X$ is totally bounded, then $A$ is bounded.
\end{proposition}

\begin{proof}
Apply total boundedness with $\varepsilon = 1$: there exist
$x_1, \ldots, x_n$ with $A \subseteq \bigcup_k B_1(x_k)$.
Fix any $p \in X$.
For each $k$, pick any $a_k \in B_1(x_k)$ (if $B_1(x_k) \cap A \neq \varnothing$;
ignore empty intersections).
Let $R = \max_k d(p, x_k) + 1$.
For any $a \in A$, there exists $k$ with $d(x_k, a) < 1$, so
$d(p, a) \le d(p, x_k) + d(x_k, a) < R + 1$.
Thus $A \subseteq B_{R+1}(p)$ and $A$ is bounded.
\end{proof}

% ---- Bounded does not imply totally bounded ----

\begin{proposition}[Bounded does not imply totally bounded]
\label{prop:bounded-not-totally-bounded}
An infinite set with the discrete metric is bounded but not totally bounded.
\end{proposition}

\begin{proof}
Let $X$ be infinite with the discrete metric $d(x,y) = 1$ for $x \neq y$.
Then $X \subseteq B_2(x_0)$ for any fixed $x_0$, so $X$ is bounded.

For $\varepsilon = 1/2$: each ball $B_{1/2}(x) = \{x\}$ contains only one point.
To cover $X$ by $1/2$-balls, we need one ball per point --- infinitely many.
So $X$ is not totally bounded.
\end{proof}

\begin{remark*}[Why the discrete metric is the key counterexample]
In the discrete metric, every point is ``far from'' every other
point in the sense that you need a separate ball to capture each one.
No finite collection of small balls can reach all of them.
This is the essence of why total boundedness fails: the space has
``infinitely many independent directions.''
\end{remark*}

% ---- Rn: bounded iff totally bounded ----

\begin{proposition}[Bounded and totally bounded coincide in $\mathbb{R}^n$]
\label{prop:totally-bounded-Rn}
A subset $A \subseteq \mathbb{R}^n$ is totally bounded if and only if
it is bounded.
\end{proposition}

\begin{proof}
($\Rightarrow$) Proposition~\ref{prop:totally-bounded-implies-bounded}.

($\Leftarrow$) Suppose $A$ is bounded, so $A \subseteq [-M, M]^n$ for
some $M > 0$.
Let $\varepsilon > 0$.
Partition $[-M, M]^n$ into a finite grid of cubes of side length
$\varepsilon / \sqrt{n}$.
The number of cubes is at most $(2M\sqrt{n}/\varepsilon + 1)^n$, which is finite.
Each cube has diameter $\sqrt{n} \cdot (\varepsilon/\sqrt{n}) = \varepsilon$,
so it is contained in an open ball of radius $\varepsilon$ centred at its midpoint.
Hence finitely many $\varepsilon$-balls cover $[-M,M]^n \supseteq A$.
\end{proof}

\begin{remark*}[Why this fails in infinite dimensions]
In $\ell^2$ (the space of square-summable sequences), the closed unit ball
$\{x : \|x\|_2 \le 1\}$ is bounded but not totally bounded.
Consider the sequence of standard basis vectors $e_n$
(all coordinates $0$ except the $n$-th which is $1$).
Each $e_n$ has norm $1$, but $\|e_m - e_n\|_2 = \sqrt{2}$ for all $m \neq n$.
So the $e_n$ are pairwise at distance $\sqrt{2}$ apart.
No $(\sqrt{2}/2)$-ball can contain more than one $e_n$,
so infinitely many balls are needed to cover $\{e_n : n \in \mathbb{N}\}
\subseteq B_1(0)$.
Hence $B_1(0)$ is not totally bounded in $\ell^2$.

This is why Heine-Borel fails in infinite dimensions: bounded $+$ closed
is not enough to force compactness.
Total boundedness is the missing condition.
\end{remark*}

% ---- Totally bounded gives Cauchy subsequences ----

\begin{proposition}[Totally bounded sets have Cauchy subsequences]
\label{prop:totally-bounded-Cauchy}
If $A \subseteq X$ is totally bounded, then every sequence in $A$ has
a Cauchy subsequence.
\end{proposition}

\begin{proof}
Let $(a_n)$ be a sequence in $A$.
We construct a Cauchy subsequence by a diagonal argument.

\medskip
\noindent\textbf{Step 1.}
Apply total boundedness with $\varepsilon = 1$:
cover $A$ with finitely many $1$-balls.
By the pigeonhole principle, one of these balls ---
call it $B_1$ --- contains $a_n$ for infinitely many $n$.
Extract a subsequence $(a_n^{(1)})$ of $(a_n)$ lying entirely in $B_1$.
All terms of $(a_n^{(1)})$ satisfy $d(a_m^{(1)}, a_k^{(1)}) < 2$.

\medskip
\noindent\textbf{Step 2.}
Apply total boundedness with $\varepsilon = 1/2$:
cover $A$ with finitely many $(1/2)$-balls.
By pigeonhole, one --- call it $B_2$ --- contains $(a_n^{(1)})_n$
for infinitely many $n$.
Extract a subsequence $(a_n^{(2)})$ of $(a_n^{(1)})$ lying in $B_1 \cap B_2$.
All terms satisfy $d(a_m^{(2)}, a_k^{(2)}) < 1$.

\medskip
\noindent\textbf{Continuing.}
At step $j$, extract a subsequence $(a_n^{(j)})$ of $(a_n^{(j-1)})$
lying in a $(1/j)$-ball, so all terms satisfy
$d(a_m^{(j)}, a_k^{(j)}) < 2/j$.

\medskip
\noindent\textbf{Diagonal subsequence.}
Define $b_j = a_j^{(j)}$ (the $j$-th term of the $j$-th subsequence).
For $m, k \ge J$, both $b_m$ and $b_k$ are terms of $(a_n^{(J)})$
(by construction: $(a_n^{(j)})$ is a subsequence of $(a_n^{(J)})$ for $j \ge J$),
so $d(b_m, b_k) < 2/J$.
Since $J$ can be taken as large as we like, $(b_j)$ is Cauchy.
\end{proof}

\begin{remark*}[This is the heart of sequential compactness]
Proposition~\ref{prop:totally-bounded-Cauchy} says: totally bounded
gives Cauchy subsequences.
Completeness then turns those Cauchy subsequences into convergent ones.
So totally bounded $+$ complete implies: every sequence has a convergent
subsequence, which is exactly sequential compactness.
This is the proof strategy behind the Generalized Heine-Borel theorem.
\end{remark*}

% ---- Generalized Heine-Borel ----

\begin{tcolorbox}[colback=thmbox, colframe=thmborder, arc=2pt,
  left=6pt, right=6pt, top=4pt, bottom=4pt,
  title={\small\textbf{Theorem (Generalized Heine-Borel --- Pugh Theorem 65)}},
  fonttitle=\small\bfseries]
\begin{theorem}[Generalized Heine-Borel]
\label{thm:gen-heine-borel}
Let $(X,d)$ be a complete metric space and $A \subseteq X$.
Then $A$ is compact if and only if $A$ is closed and totally bounded.
\end{theorem}
\end{tcolorbox}

\begin{remark*}[Proof strategy --- deferred to the compactness chapter]
($\Rightarrow$) Every compact set is closed (as a complete metric space
in its own right) and totally bounded (compactness gives a finite subcover
for any open cover by $\varepsilon$-balls).

($\Leftarrow$) Let $(a_n)$ be any sequence in $A$.
By Proposition~\ref{prop:totally-bounded-Cauchy}, it has a Cauchy
subsequence $(b_n)$.
Since $X$ is complete, $(b_n)$ converges to some $p \in X$.
Since $A$ is closed, $p \in A$.
So every sequence in $A$ has a subsequence converging in $A$,
which is sequential compactness, hence compactness.
The full proof (including the equivalence of sequential and open-cover compactness)
appears in the compactness chapter.
\end{remark*}

\begin{remark*}[Comparison with classical Heine-Borel]
In $\mathbb{R}^n$, Proposition~\ref{prop:totally-bounded-Rn} shows
that totally bounded is the same as bounded.
So in $\mathbb{R}^n$, the Generalized Heine-Borel theorem reads:
\emph{compact $\iff$ closed and bounded.}
This is the classical Heine-Borel theorem.

The generalized version reveals what the classical theorem is really saying:
``bounded'' in $\mathbb{R}^n$ happens to coincide with
``totally bounded,'' but it is total boundedness --- not mere boundedness ---
that is the correct hypothesis for compactness in general metric spaces.
\end{remark*}

\begin{example}[Compact vs.\ merely closed-and-bounded in $C([0,1])$]
\label{ex:ell2-not-totally-bounded}
The closed unit ball $B = \{f \in C([0,1]) : \|f\|_\infty \le 1\}$ in
$C([0,1])$ with the uniform metric is closed and bounded.

However, $B$ is \emph{not} totally bounded.
Consider the functions $f_n(x) = x^n$.
Each $\|f_n\|_\infty = 1$, so $f_n \in B$.
One can check that $\|f_m - f_n\|_\infty = 1$ for all $m \neq n$
(since $f_m - f_n$ achieves the value $1/2$ near $x = 1$ for appropriate
choice of $m, n$).
So no finite collection of $(1/4)$-balls can cover $\{f_n\}$,
and $B$ is not totally bounded, hence not compact.

This confirms: in infinite-dimensional spaces, closed and bounded is not
enough.
\end{example}
